% !TeX root = ../papers/nc-p3.tex
\section{階数15の反対切断の核線と添字保持同型}
\label{sec:rank-fifteen-opposite-kernels}
\begin{theorem}\label{thm:rank-fifteen-opposite-kernels}
定理\ref{thm:binary-opposite-determinant}の残余因子の一般点で、
反対切断の左右の核線はどちらも非退化な四元双線型形式で生成される。
\end{theorem}
\begin{proof}
同定理の有理標本で、滑らかな15次元枝の残余因子内に形式曲線を取る。
二交換のそれぞれで符号を変える座標を $a,b$ とし、その一次方向を同定理の $A,B$ に選ぶ。
$a=b=t$ とおき、残りの不変座標を $t^2$ の級数に取れる。
実際、交叉上でも $D$ の微分は非零なので、一つの不変座標を $D=0$ で消去できる。
曲線の二次までの不変座標は、15次元枝の式と $D=0$ の微分を満たす範囲で選べる。
ここで $a,b$ は文字基底の軌道も含め同変な滑らかな座標を取り直してよい。

$w(t)$ の一次項を $A+B$、二次項を $X+Z$、三次項を $A_3+B_3$ とし、
四次項の交代・交代部分を $Z_4$ と書く。
$X,Z,A_3,B_3$ の指標は順に $++,--,+-,-+$ である。
対向する二組をそれぞれ対称・交代に分けた基底で、実際の反対切断は
$\operatorname{diag}(I_{10},tI_6)M(t)\operatorname{diag}(I_{10},tI_6)$ である。
$M(t)=M_0+t^2M_2+O(t^4)$ の四ブロックは、順に
$(w,A,B,Z)$ と $(X,A_3,B_3,Z_4)$ の反対切断の対応部分から取る。

係数の決め方を記す。第一隣接切断の核列を $N$、一般逆を $J$、
増分を $E$ とし、Schur方程式
\[
 N^t(E-EJE+EJEJE-EJEJEJE+\cdots)N=0
\]
の指定次数を取る。もう一つの隣接切断も課すには、
二変数 $a,b$ の各指標の項を別々に足す。これは一方の交換が移す切断にも
同じ式を課すことに等しい。二次で $X,Z$ を解き、三次の混合障害が消えるよう
$X$ に対称接を加える。定理\ref{thm:binary-opposite-determinant}の滑らかな枝の式により、
二つの三次の障害は同じ一条件である。$A_3,B_3$ の自由接成分は基底と混合座標の
三次補正で選べ、$Z_4$ は第四次の $--$ 方程式から一意に決まる。
最後に $M_0$ の左右の核生成元 $l,r$ に対し $l^tM_2r=0$ を課す。
$X$ に交叉に接する対称方向を加えたときのこの式の微分は前定理の非零微分である。
これで残余因子内の曲線の必要な係数が決まる。以降の係数は上の陰関数消去で延長できる。

\path{computations/verify_rank_fifteen_opposite_kernels.py} はこれらの線型系を素数1009で
解き、全方程式への代入も検査する。$l=(C,0)$ の対称形式 $C$ は第四行・列が零で
左上 $3\times3$ ブロック $C_3$ が可逆であり、$r$ の交代部分は階数4である。
従って実際の右核の生成元を $t$ 倍した極限は可逆な交代形式である。
左側では $l_2^t M_0=-l^tM_2$ を解くと、実際の核形式は
$C+tK+t^2L+O(t^3)$ となる。$K,L$ は $l_2$ の交代・対称部分で、
$K_{44}=0$ である。左上ブロックに関するSchur補の $t^2$ の係数は
\[
 L_{44}-K_{4,1:3}C_3^{-1}K_{1:3,4}\equiv 22\pmod {1009}.
\]
$l_2$ に $l$ を加える自由度はこの係数を変えない。
前定理で標数零とこの素数で隣接Jacobi階数、混合障害の階数および
残余因子の横断性が一致しており、全ての線型系は同じ可逆小行列で持ち上がる。
$M_0$ の階数15の上界・錐の式も標数零で既知である。従って左右の非零小行列は
標数零でも非零である。形式曲線の一般点での非消滅と開性から主張が従う。
\end{proof}

\begin{proposition}\label{prop:opposite-adjoint-equivalence}
任意の反対階数15の四窓テンソルで、核双線型形式が両方とも非退化とする。
四文字空間を順に $A,B,A^*,B^*$ と同定すると、この部分の添字保持同型は
射影線型同型 $[f]\colon\mathfrak{sl}(A)\longrightarrow\mathfrak{sl}(B)$ の
独立な左右の内部共役にちょうど一致する。一添字移動の式は、トレース随伴を
$f^\dagger$、双対空間に値を取る転置を $\tau$ として $[\tau_A f^\dagger]$ であり、
二添字移動は $[\tau_B f\tau_A^{-1}]$ である。
\end{proposition}
\begin{proof}
それぞれの核線を標準評価に取る。その縮約が零だからテンソルは
$\mathfrak{sl}(A)\otimes\mathfrak{sl}(B)$ に属し、非退化なトレース対で
線型写像 $f$ を与える。反対階数15がその可逆性に等しい。
標準評価の線を保つ基底変更は一方の文字の基底変更と反対文字での逆転置のスカラー倍に限られる。
これが $f\mapsto \lambda\operatorname{Ad}_h f\operatorname{Ad}_g^{-1}$ を与え、
逆に全てのそのような変換が実現する。テンソル $\sum_i P_i\otimes Q_i$ は
一添字移動で $\sum_i Q_i\otimes P_i^t$ となるので移動の式が従う。
従って各移動と元の写像の内部共役二重軌道の一致が、その移動が周期となる必要十分条件である。
窓の隣接階数やFrobenius条件には、$f$ から戻した同一テンソルを用いる。
\end{proof}

四文字に体積形式を選ぶと、反対切断の余因子行列の四乗を
各文字の交代体積で縮約して次数60の $\mathrm{SL}_4^4$ 不変式を得る。
反対階数15では余因子が核生成元の単純積 $x\otimes y$ であるから、この値は
非零定数倍を除き $\det(x)\det(y)$ である。従って命題の開部分はグラフの次元によらず半安定である。
半安定商における閉軌道代表でもこの不変式は非零であり、階数15以下は閉条件であるため、
同じ正規化でその代表を調べられる。
